\section{Argyres-Douglas 共形场论}\label{sec:AD}\index{ArgyresDouglasCFT@Argyres-Douglas 共形场论|hyperpage}\index{ArgyresDouglasPoint@Argyres-Douglas 点|hyperpage}\index{ScalingDimension@标度维数|hyperpage}
本节回到第 \ref{sec:nf_1_curve} 小节末尾的观察：$N_f=1$ 理论的库仑支上存在一个高度奇异的点。我们将研究该点的物理及其推广。

\subsection{$N_f=1$ 理论与最简单的 Argyres-Douglas 共形场论}
再次回到含有 $N_f=1$ 个味的 $\SU(2)$ 规范理论的曲线：
\begin{equation}
\Sigma:\quad\frac{2\Lambda(x-\mu)}z+\Lambda^2z=x^2-u.\label{nf1curveagain}
\end{equation}
对于一般的 $\Lambda$ 和 $\mu$，$u$ 平面上有三个奇点。正如第 \ref{sec:nf_1_curve} 小节末尾所见，当取 $\mu=-\frac32\Lambda$ 时，两个奇点在 $u=3\Lambda^2$ 处碰撞，见图 \ref{fig:nf_1_adpoint}。
\begin{figure}[h]
\[
\inc{ADnf1X}
\]
\caption{$N_f=1$ 理论的 Argyres-Douglas 点。\label{fig:nf_1_adpoint}}
\end{figure}
当 $u\sim3\Lambda^2$ 时，$x(z)$ 的三个分支点在 $z=-1$ 处碰撞。此时，定义 $a$ 和 $a_D$ 的 $A$ 闭链与 $B$ 闭链，都可以取为环绕 $z=0$ 的小环路。这保证了 $a$ 和 $a_D$ 都很小。因此，我们同时具有很轻的电荷粒子和磁荷粒子。库仑支上的这种点称为 Argyres-Douglas 点。文献 \cite{Argyres:1995jj} 首先在纯 $\SU(3)$ 理论中识别出了这样的点，文献 \cite{Argyres:1995xn} 又将其推广到含有味的 $\SU(2)$ 理论。

% [Tutor Note: p.109 原文 "small loops around $z=0$" 中的 $z=0$ 应为 $z=-1$：同一段前一句刚说明三个分支点在 $z=-1$ 处碰撞，围绕这些分支点的小环路应取在 $z=-1$ 附近（v3 版已改作 $z=-1$，可作旁证）。正文照原文排 $z=0$。]

环绕 $u=3\Lambda^2$ 的单值变换 $M_{AD1}$ 可以用多种方法求得。一种方法是将 $N_f=1$ 理论中两个相碰奇点的单值变换相乘。另一种方法是追踪三个分支点的运动。令 $u=3\Lambda^2+\delta u$，可得三个分支点位于 $z+1\propto\delta u^{1/3}$ 处。这就确定了闭链如何被映射，从而得到单值变换。无论使用哪种方法，都有
\begin{equation}
M_{AD1}\sim\begin{pmatrix}
1&1\\
-1&0
\end{pmatrix}.
\end{equation}
$M_{AD1}$ 对低能耦合的变换为
\begin{equation}
\tau\mapsto\tau'=\frac1{1-\tau}.
\end{equation}
注意，$\tau=e^{\pi i/3}$ 是这一变换的不动点；通过显式计算，可以检验 $\tau-e^{\pi i/3}\propto\delta u^{1/3}$。我们发现，耦合被固定在这一强耦合值上。

人们认为，低能极限是共形的。为了分离出这一极限中的物理，令
\begin{equation}
\begin{aligned}
z&=-1+\tilde\delta z,&
x&=-\Lambda+\Lambda\delta x,&
u&=3\Lambda^2+\Lambda^2\tilde\delta u,&
\mu&=-\frac32\Lambda+\Lambda\tilde\delta\mu.
\end{aligned}
\end{equation}
并假定所有以 $\delta$ 为前缀的变量都很小。事实表明，进一步引入如下重定义更加有用：
\begin{equation}
\tilde\delta z=\delta z-(\delta z)^2,\quad
\tilde\delta u=\delta u-2\delta\mu-2\delta x+(\delta z)^2,\quad
\tilde\delta\mu=\delta\mu+\delta x.\label{ADredefinition1}
\end{equation}

现在，将关系 \eqref{ADredefinition1} 代入 $N_f=1$ 曲线 \eqref{nf1curveagain}。展开到 $\delta z$ 的三阶，得到如下形式的曲线：
\begin{equation}
(\delta x)^2-\delta u=-(\delta z)^3+2\delta\mu\delta z.\label{nf1ADcurve}
\end{equation}
微分为
\begin{equation}
\lambda=x\frac{dz}z\sim\delta x d\delta z.
\end{equation}
由于积分 $\int\lambda$ 给出 BPS 粒子的质量，$\lambda$ 本身应具有标度维数 $1$。关系 \eqref{nf1ADcurve} 意味着，标度维数 $[\delta x]$ 与 $[\delta z]$ 应满足
\begin{equation}
[\delta x]:[\delta z]=3:2.
\end{equation}
这就固定了所有相关变量的标度维数：
\begin{equation}
[\delta x]=\frac35,\qquad[\delta x]=\frac25,\qquad
[\delta u]=\frac65,\qquad[\delta\mu]=\frac45.
\end{equation}
% [Tutor Note: p.110 式 (10.1.9) 原文第二个 $[\delta x]$ 应为 $[\delta z]$：上一式给出 $[\delta x]:[\delta z]=3:2$，下面两式依次使用 $[\delta x]=3/5$、$[\delta z]=2/5$、$[\delta u]=6/5$、$[\delta\mu]=4/5$（v3 版已改作 $[\delta z]$，可作旁证）。正文照原文排。]
注意，算符 $u=\tr\Phi^2/2$ 的质量量纲，等价地说，它在紫外的标度维数，原本是 $2$。我们发现，反常维数是 $1$ 的量级，使 $[\delta u]$ 显著减小。

\begin{figure}[h]
\[
\inc{AD1}
\]
\caption{Argyres-Douglas 理论 $AD_{N_f=1}(\SU(2))$。\label{fig:H1}}
\end{figure}
在取极限 $\delta u\to0$ 时，我们是在放大 $u$ 平面上 $u=3\Lambda^2$ 附近的邻域。在这一极限中，可以认为低能理论由一个仅在 $\delta u=0$ 处具有奇点的理论描述，如图 \ref{fig:H1} 左侧所示。我们将得到的理论称为 Argyres-Douglas 共形场论 $AD_{N_f=1}(\SU(2))$\footnote{遗憾的是，文献中对于 Argyres-Douglas 理论还没有公认的统一命名体系。本讲义试图提供一种命名体系，它可能比文献中的命名更加繁琐，但在区分各种构造时更加明确。}。

从六维观点重新考察这一取极限的过程。首先，将原来的 $N_f=1$ 曲线写成 $\lambda^2-\phi(z)=0$ 的形式。回顾第 \ref{sec:less} 小节的研究以及图 \ref{fig:lessG}，$\phi(z)$ 有一个三阶极点和一个四阶极点。此外，在一般位置上还有 $\phi(z)$ 的三个分支点。假定仔细调节参数，使 $\phi(z)$ 的两个极点碰撞：
\begin{equation}
\phi(z)\sim(\frac{P_3(z)}{(z-\epsilon)^3}+\frac{P_4(z)}{z^4})dz^2
=\frac{P_7(z)}{(z-\epsilon)^3z^4}dz^2
\to\frac{Q_7(z)}{z^7}dz^2
\end{equation}
其中，在这一阶段，$P_d$、$Q_d$ 是一般的 $d$ 次多项式。最终只剩下一个具有七阶极点的奇点，如图 \ref{fig:H1} 右侧所示。为了在 $z=\infty$ 处没有奇点，$Q_7(z)$ 实际上应当是三次多项式：
\begin{equation}
\lambda^2=\frac{c+c'z+\mu z^2+uz^3}{z^7}dz^2.
\end{equation}
通过坐标变换 $z\to z/(az-b)$，可以令 $c=1$、$c'=0$。于是有
\begin{equation}
\lambda^2=\frac{1+\mu z^2+uz^3}{z^7}dz^2.
\end{equation}
由于左侧的标度维数为 $2$，可得 $[z]=-2/5$，从而有
\begin{equation}
[\mu]=\frac45,\qquad[u]=\frac65\label{muu}
\end{equation}
这与前面的结果一致。注意，变量 $z$ 是辅助变量，因此没有理由要求它的维数与前面相同。

一般而言，对于任意共形场论，任意动力学标量算符 $O$ 都应具有大于或等于 $1$ 的标度维数：
\begin{equation}
[O]\ge1,\label{generalCFTfeatures}
\end{equation}
而且只有当 $O$ 描述一个自由、退耦的标量玻色子时，等号才能成立。因此，具有 $[u]=6/5$ 的算符 $u$ 是理论 $AD_{N_f=1}(\SU(2))$ 中的真正算符。$\mu$ 则在如下意义下被视为与 $u$ 共轭的参数。在 $\cN{=}2$ 理论中，可以考虑预势的形变
\begin{equation}
\int d^4\theta F\to\int d^4\theta(F+mO).
\end{equation}
这里，$d^4\theta$ 是第 \ref{sec:lel} 小节末尾简要提到的手征 $\cN{=}2$ 超空间积分，$O$ 是一个算符，$m$ 是乘在它前面的参数。因此，在 $\cN{=}2$ 超共形理论中，组合 $mO$ 需要具有标度维数 $2$。于是应有
\begin{equation}
[m]+[O]=2.
\end{equation}

可以看到，$\mu$ 与 $u$ 这一对满足此条件，见式 \eqref{muu}。因此，我们将 $\mu$ 视为对应于算符 $u$ 的形变参数。

\subsection{来自 $N_f=2$ 理论的 Argyres-Douglas 共形场论}\index{MassParameter@质量参数|hyperpage}
考虑 $N_f=2$ 理论的曲线，
\begin{equation}
\frac{2\Lambda(x-\mu)}z+2\Lambda(x-\mu)z=x^2-u\label{firstform}
\end{equation}
其中令两个味的质量相同。也可以使用另一种形式的曲线
\begin{equation}
\frac{(x-\mu)^2}z+4\Lambda^2z=x^2-u.\label{secondform}
\end{equation}
对于一般的 $\mu$，其模空间已经画在图 \ref{fig:nf_2_moduli} 中。

当 $\mu=0$ 时，两个没有连接希格斯支的奇点发生碰撞。现在，转而调节参数 $\mu$，使带有希格斯支的奇点与一个不带希格斯支的奇点碰撞，见图 \ref{fig:nf_2_AD_plane}。为了明确起见，使用曲线的第一种形式。此时，碰撞发生在 $\mu=2\Lambda$、$u=4\Lambda^2$ 时。四个分支点随后在 $z=1$ 处碰撞。
\begin{figure}[h]
\[
\inc{ADnf2X}
\]
\caption{$N_f=2$ 理论的 Argyres-Douglas 点。\label{fig:nf_2_AD_plane}}
\end{figure}
我们发现，环绕所得到奇点的单值变换为
\begin{equation}
M_{AD2}\sim\begin{pmatrix}
0&1\\
-1&0
\end{pmatrix},
\end{equation}
它对耦合的作用为
\begin{equation}
\tau\to\tau'=-\frac1\tau.
\end{equation}
强耦合值 $\tau=i$ 是这一变换的不动点，当令 $u\to4\Lambda^2$ 时，低能耦合趋近于这一值。

与前面一样展开各变量，
\begin{equation}
\begin{aligned}
z&=1+\delta z,&
x&=2\Lambda+\delta x,&
u&=4\Lambda^2+\delta u,&
\mu&=2\Lambda+\delta\mu,
\end{aligned}
\end{equation}
可得极限中的曲线为
\begin{equation}
(\delta x)^2+\delta u=(\delta z)^4+\delta\mu\delta z^2+\Delta\mu\delta z\label{nf2ADcurve}
\end{equation}
微分为 $\lambda\sim\delta x d\delta z$。这里，重新引入了两个二重态超多重态的裸质量 $\mu_1$、$\mu_2$ 之间的微小差值 $\Delta\mu=\mu_1-\mu_2$。要求 $\lambda$ 的标度维数为 $1$，便有
\begin{equation}
[\delta x]=\frac23,\quad[\delta z]=\frac13.
\end{equation}
于是得到
\begin{equation}
[\delta u]=\frac43,\quad[\delta\mu]=\frac23,\quad[\Delta\mu]=1.\label{scalingdimH2}
\end{equation}
再次可见，$[\delta u]+[\delta\mu]=2$，因此 $\delta\mu$ 是对应于算符 $\delta u$ 的形变参数。$\Delta\mu$ 是非阿贝尔味对称性 $\SU(2)_F$ 的质量参数。一般而言，在共形理论中，非阿贝尔味对称流 $J^a$ 应具有标度维数 $3$。$\cN{=}2$ 超对称性将它与质量项联系起来；对于拉格朗日理论，这个质量项就是熟悉的 $Q\tilde Q$，具有标度维数 $2$。因此，$\cN{=}2$ 超共形理论的非阿贝尔质量参数始终应具有标度维数 $1$。我们对 $[\Delta\mu]$ 的计算与这一一般论证一致。将得到的理论称为 $AD_{N_f=2}(\SU(2))$。

\begin{figure}[h]
\[
\inc{AD2}
\]
\caption{Argyres-Douglas 理论 $AD_{N_f=2}(\SU(2))$。\label{fig:H2}}
\end{figure}
从六维观点研究 $N_f=2$ 理论的取极限过程。在取极限之前，第一种形式 \eqref{firstform} 的曲线为 $\lambda^2=\phi(z)$，其中 $\phi(z)$ 具有两个四阶极点。使它们碰撞，最终得到一个八阶奇点。与前面的分析一样，可知极限中的曲线为
\begin{equation}
\lambda^2=\frac{1+\mu z^2+\Delta\mu z^3+uz^4}{z^8}dz^2.
\end{equation}
很容易看出 $[z]=-1/3$。因此，得到的标度维数与式 \eqref{scalingdimH2} 相同。

第二种形式 \eqref{secondform} 的曲线在写成 $\lambda^2=\phi(z)$ 时，具有两个二阶极点和一个三阶极点。在两个二阶极点处，$xdz/z$ 的留数分别为 $\pm(\mu_1+\mu_2)/2$ 和 $\pm(\mu_1-\mu_2)/2$。将留数为 $\pm(\mu_1+\mu_2)/2$ 的二阶极点与三阶极点碰撞，形成一个五阶极点。最终得到的 $\phi(z)$ 具有一个五阶极点，不妨设在 $z=0$，以及另一个留数为 $\pm(\mu_1-\mu_2)/2$ 的二阶极点，见图 \ref{fig:H2}。极限中的曲线也很容易求得：
\begin{equation}
\lambda^2=\frac{1+\delta\mu z+\delta uz^2+(\frac{\mu_1-\mu_2}2)^2z^3}{z^5}dz^2
\end{equation}
最后一个系数由 $z=\infty$ 处的条件固定。要求 $\lambda$ 具有标度维数 $1$，可得 $[z]=-2/3$，以及
\begin{equation}
[\delta\mu]=\frac23,\quad[\delta u]=\frac43.
\end{equation}
得到相同的答案，令人安心。

\subsection{来自 $N_f=3$ 理论的 Argyres-Douglas 共形场论}
可以用完全相同的方法求出 $N_f=3$ 理论的特殊极限。从曲线 \eqref{nf3curve!} 出发，
\begin{equation}
\frac{(x-\mu-\Lambda)^2}z+2\Lambda(x-\mu-\Lambda)z=x^2-u
\end{equation}
令三个味具有相同质量。在 $u$ 平面上，有一个带有希格斯支的奇点，以及两个不带希格斯支的奇点。调节 $\mu$，使带有希格斯支的奇点与另一个不带希格斯支的奇点碰撞，并使它们的单值变换不对易。见图 \ref{fig:nf_3_AD_uplane}。

\begin{figure}[h]
\[
\inc{ADnf3}
\]
\caption{$N_f=3$ 理论的 Argyres-Douglas 点。\label{fig:nf_3_AD_uplane}}
\end{figure}
环绕所得到奇点的单值变换为
\begin{equation}
M_{AD3}=\begin{pmatrix}
0&1\\
-1&-1
\end{pmatrix}
\end{equation}
它对耦合的作用为
\begin{equation}
\tau\mapsto\tau'=\frac{-\tau+1}{-\tau}.
\end{equation}
不动点位于 $\tau=e^{\pi i/3}$。

\begin{figure}[h]
\[
\inc{AD3}
\]
\caption{Argyres-Douglas 理论 $AD_{N_f=3}(\SU(2))$。\label{fig:H3}}
\end{figure}
在六维描述中，原来有两个二阶极点和一个四阶极点。使一个二阶极点与一个四阶极点碰撞，最终得到一个六阶极点。此时曲线为
\begin{equation}
\lambda^2=\frac{1+\delta\mu z+\mu'z+\delta uz^2+(\frac{\mu_1-\mu_2}2)^2z^3}{z^6}dz^2
\end{equation}
微分 $\lambda$ 的标度维数为 $1$。于是 $[z]=-1/2$，并且有
\begin{equation}
[\delta\mu]=\frac12,\quad[\delta u]=\frac32,\quad[\mu']=1,
[\Delta\mu]=1,
\end{equation}
其中定义了 $\Delta\mu=\mu_1-\mu_2$。两个参数 $\mu'$ 和 $\Delta\mu$ 的标度维数为 $1$，我们将它们与 $\SU(3)$ 味对称性相联系的质量参数对应起来。同样再次得到 $[\delta\mu]+[\delta u]=2$。将所得理论称为 $AD_{N_f=3}(\SU(2))$。

\subsection{秩为一的理论的总结}\label{sec:rank-1SCFT}\index{RankOneTheory@秩为一的理论|hyperpage}\index{KodairaClassification@Kodaira 分类|hyperpage}
\subsubsection{来自含味 $\SU(2)$ 理论的 Argyres-Douglas 共形场论}
迄今为止，我们研究了由含有 $N_f=1,2,3$ 个味的 $\SU(2)$ 规范理论的特殊极限得到的 Argyres-Douglas 共形场论。这些以及其他相关共形场论的数据汇总在表 \ref{tab:rank1sing} 中。Argyres-Douglas 共形场论位于表的前三行。第四行对应含有 $N_f=4$ 个无质量味的 $\SU(2)$ 理论。接下来的两行对应略有不同的理论类别。具体而言，如果考虑味数大于 $4$ 的 $\SU(2)$ 理论，或含有非零数目带荷超多重态的 $\U(1)$ 理论，那么它们是红外自由的，见式 \eqref{su2running} 和式 \eqref{general_u1_running}。它们曾反复出现，用来描述 $u$ 平面上奇点附近的局部行为。

表 \ref{tab:rank1sing} 中还列出了希格斯支的维数。我们来简要回顾如何得到这些维数。已知 $AD_{N_f=1}(\SU(2))$ 没有希格斯支，因为它的母理论，即含有 $N_f=1$ 个味的 $\SU(2)$ 理论，也没有希格斯支。对于 $AD_{N_f=2}(\SU(2))$，考虑带有 $\U(1)$ 质量项的 $N_f=2$ $\SU(2)$ 理论。此时希格斯支为 $\bC^2/\bZ_2$，其四元数维数为 $1$。对于 $AD_{N_f=3}(\SU(2))$，考虑 $N_f=3$ 的 $\SU(2)$ 理论。加入 $\U(1)$ 质量项后，可以通过研究弱耦合理论中的点 $u=\mu_1^2=\mu_2^2=\mu_3^2$ 来求得希格斯支。该处的物理就是含有三个带荷超多重态的 $\U(1)$ 理论，希格斯支的四元数维数为 $2$。对于含有 $N_f\ge4$ 个味的自由 $\SU(2)$ 理论，四元数维数就是 $2N_f-\dim\SU(2)$；类似地，对于含有 $N$ 个味的 $\U(1)$ 理论，它就是 $N-\dim\U(1)$。一个有趣的特点是，对于前六行，有
\begin{equation}
\dim_\bH(\text{希格斯支})=h^\vee(\text{味对称性})-1\label{hf}
\end{equation}
其中 $h^\vee(G)$ 是对偶 Coxeter 数，它也等于 $G$ 的伴随表示对单圈跑动的贡献 $C(\text{adj})$。这些理论只有一个库仑支模，而库仑支上一般点处的低能理论只是一个自由 $\U(1)$ 理论。这样的理论称为秩为一的理论。

\begin{table}
\[
\begin{array}{c|cccccccccccc}
\text{名称}&\text{单值变换}&\text{味对称性}&[u]&\#&\dim_\bH(\text{希格斯支})\\
\hline
AD_{N_f=1}(\SU(2))&\begin{pmatrix}1&1\\-1&0\end{pmatrix}&&6/5&2&\\
AD_{N_f=2}(\SU(2))&\begin{pmatrix}0&1\\-1&0\end{pmatrix}&\SU(2)&4/3&3&1\\
AD_{N_f=3}(\SU(2))&\begin{pmatrix}0&1\\-1&-1\end{pmatrix}&\SU(3)&3/2&4&2\\
\SU(2)\ N_f=4&\begin{pmatrix}-1&0\\0&-1\end{pmatrix}&\SO(8)&2&6&5\\
\hline
\SU(2)\ N_f>4&\begin{pmatrix}-1&4-N_f\\0&-1\end{pmatrix}&\SO(2N_f)&&N_f+2&2N_f-3\\
\text{含 $N$ 个味的}\ \U(1)&\begin{pmatrix}1&N\\0&1\end{pmatrix}&\SU(N)&&N&N-1\\
\hline
MN(E_6)&\begin{pmatrix}-1&-1\\1&0\end{pmatrix}&E_6&3&8&11\\
MN(E_7)&\begin{pmatrix}0&-1\\1&0\end{pmatrix}&E_7&4&9&17\\
MN(E_8)&\begin{pmatrix}0&-1\\1&1\end{pmatrix}&E_8&6&10&29
\end{array}
\]
\caption{各种秩为一的共形场论的数据。$\#$ 是在 $u=0$ 处碰撞的奇点数目，$\dim_\bH\text{希格斯支}$ 是希格斯支的四元数维数，即实维数除以 $4$。\label{tab:rank1sing}}
\end{table}

\begin{table}
\[
\begin{array}{c||c|c|c|c|c}
\Gamma&\bZ_n&\widehat{\mathcal D}_{n-2}&\widehat{\mathcal T}&\widehat{\mathcal O}&\widehat{\mathcal I}\\
\hline
G_\Gamma&\SU(n)&\SO(2n)&E_6&E_7&E_8
\end{array}
\]
\caption{$\SU(2)$ 的有限子群 $\Gamma$ 与单链李群 $G_\Gamma$。这里，$\mathcal D_n$ 是作用于正 $n$ 边形的二面体群，$\mathcal T$、$\mathcal O$、$\mathcal I$ 分别是四面体群、八面体群和二十面体群，上方的帽号表示从 $\SO(3)$ 到 $\SU(2)$ 的提升。例如，得到的群 $\widehat{\mathcal T}$ 称为二元四面体群。\label{tab:ADE}}
\end{table}

\subsubsection{Minahan-Nemeschansky 的例外理论}\label{sec:F}\index{MinahanNemeschanskyTheory@Minahan-Nemeschansky 理论|hyperpage}\index{FTheory@F 理论|hyperpage}\index{ExceptionalFlavorSymmetry@例外味对称性|hyperpage}
我们尚未讨论剩下三行列出的理论。引出这些理论的一种方式，是借助 Kodaira 的一个经典数学结果。在 $u$ 平面上的给定点，有紫外曲线 $C$ 和 Seiberg-Witten 曲线 $\Sigma$。曲线 $\Sigma$ 是一个环面，其形状由复结构 $\tau$ 参数化，而 $\tau$ 全纯地依赖于 $u$。因此，在以 $u$ 为坐标的复平面上，有一个环面纤维化。$u$ 平面连同纤维 $\Sigma$ 构成一个复二维空间 $X$。

Kodaira 对这种纤维化可能出现的奇点类型作了分类，而表 \ref{tab:rank1sing} 的前六行，正是该分类中一部分内容的完整复现。当然，术语有所不同，因为他是数学家，而我们研究的是 $\cN{=}2$ 规范理论。除前六行外，Kodaira 的分类还有三行，这促使人们认为，应当存在与之对应的另外三个理论。Minahan 和 Nemeschansky 在文献 \cite{Minahan:1996fg,Minahan:1996cj} 中首次构造了这些理论的 Seiberg-Witten 曲线。

用数学语言来说，环面纤维化中的奇点在复二维总空间 $X$ 中产生一个奇点。它局部具有 $\bC^2/\Gamma$ 的形式，其中 $\Gamma$ 是 $\SU(2)$ 的有限子群。它们具有自然的 ADE 分类，而且可以对应一个李群 $G_\Gamma$，见表 \ref{tab:ADE}。

数学家纯粹从数学上将这个群 $G_\Gamma$ 与环面纤维化联系起来，而我们看到，至少对于前六种情形，它们恰好就是规范理论的味对称性。数学家将例外群 $E_{6,7,8}$ 与最后三种情形联系起来。因此，一个很有吸引力的想法是：对应最后三行的假想理论，应以这些例外群为味对称性。根据联系味对称性与希格斯支维数的性质 \eqref{hf}，我们也很自然地想猜测这些理论的希格斯支维数。将这些共形场论分别称为 $MN(E_6)$、$MN(E_7)$ 和 $MN(E_8)$。

注意，在经典的拉格朗日 $\cN{=}2$ 理论中，很难出现例外味对称性。我们已经知道拉格朗日量的一般形式：作为 $\cN{=}1$ 理论，其超势被迫具有形式
\begin{equation}
\sum_i\int Q_i\Phi\tilde Q_i,
\end{equation}
而且可以显式检验，在紫外可见的味对称性是 $\SU$、$\SO$ 和 $\Sp$ 群的乘积 \cite{McOrist:2013bga}。因此，若要出现例外对称性，就必须通过强耦合效应产生。读者读到本讲义第 \ref{sec:application} 小节时，将会看到这在场论框架中究竟如何发生。

统一构造表中各理论的另一种方式，是使用 IIB 型弦理论及其非微扰版本 F 理论。对于含有四个味的 $\SU(2)$ 理论，这一方法始于文献 \cite{Banks:1996nj}。一般情形可参见例如文献 \cite{DeWolfe:1998bi}。采用这一方法，可以最为统一地构造这些秩为一的理论的 Seiberg-Witten 曲线，见例如文献 \cite{Noguchi:1999xq}。

IIB 型弦理论是十维的，其中有称为 $7$ 膜和 $3$ 膜的对象，$p$ 膜沿着 $p$ 个空间方向和一个时间方向延展。设时空具有形式
\begin{equation}
\bR^{1,3}\times\bR^2\times\bR^4.
\end{equation}
在子空间
\begin{equation}
\bR^{1,3}\times\{0\}\times\bR^4\label{7b}
\end{equation}
上放置一张 $7$ 膜，并在子空间
\begin{equation}
\bR^{1,3}\times\{u\}\times\{0\}.
\end{equation}
上放置一张 D3 膜。F 理论中有各种类型的 $7$ 膜，对应于 Kodaira 的分类。它们都可以由若干张最简单的 $7$ 膜，即所谓的 $(p,q)$ $7$ 膜得到：先让它们沿 $\bR^2$ 方向彼此分开，再将它们塌缩到同一点。此时，D3 膜上的低能理论就给出相应的 $\cN{=}2$ 理论。

由于存在张力，一张 $(p,q)$ $7$ 膜会产生 $\pi/6$ 的亏角。当 $n$ 张 $(p,q)$ $7$ 膜塌缩到一点时，剩下的角度为 $2\pi$ 的 $1-n/12$。由此可以计算出 $u$ 的标度维数为
\begin{equation}
u=\frac{12}{12-n},
\end{equation}
这解释了表 \ref{tab:rank1sing} 中一个有趣的规律。由若干张 $(p,q)$ $7$ 膜塌缩而得到的这些 $7$ 膜，其八维世界体积上存在规范对称性 $F$。从这一观点看，$7$ 膜上的规范对称性 $F$ 成为了味对称性。D3 膜可以作为式 \eqref{7b} 中内部 $\bR^4$ 方向上的瞬子，被这张 $7$ 膜吸收。于是，希格斯支应当由群 $F$ 的单瞬子模空间给出。群 $F$ 的 $k$ 瞬子模空间具有四元数维数 $kh^\vee(F)-1$，这就解释了关系 \eqref{hf}。

\subsubsection{较新的秩为一的理论}
迄今为止，我们看到，秩为一的理论的结构与表 \ref{tab:rank1sing} 中列出的 Kodaira 分类紧密对应。在继续之前，应当提到，还有更多秩为一的理论；它们最初是在文献 \cite{Argyres:2010py} 对各种规范理论的 S 对偶的分析中发现的。文献 \cite{Chacaltana:2012ch} 第 7 节从六维构造的观点综述了它们的性质。

\subsection{更一般的 Argyres-Douglas 共形场论：$X_N$ 与 $Y_N$}\label{sec:generalAD}\index{XNTheory@$X_N$ 理论|hyperpage}\index{YNTheory@$Y_N$ 理论|hyperpage}\index{WildGaugeTheory@野性规范理论|hyperpage}\index{CurrentTwoPointFunction@流的两点函数|hyperpage}
\begin{figure}[h]
\[
\inc{higherAD}
\]
\caption{一个更高阶的 Argyres-Douglas 理论。\label{fig:higherAD}}
\end{figure}

现在转换一下方向，考虑由规范群形如 $\SU(2)^n$ 的更复杂规范理论得到的其他 Argyres-Douglas 共形场论。例如，考虑第 \ref{sec:less} 小节末尾研究的、规范群为 $\SU(2)^4$ 的相当复杂的理论。进行与含有 $N_f=1,2,3$ 个味的 $\SU(2)$ 理论中相同的取极限过程，就可以考虑由 $\lambda^2-\phi(z)=0$ 描述的理论，其中 $\phi(z)$ 可以具有阶数很高的极点。图 \ref{fig:higherAD} 中的例子，或者只有一个十三阶极点，或者有一个九阶极点和一个十一阶极点。它们描述复杂的四维 $\cN{=}2$ 超对称共形场论。

\begin{figure}[h]
\[
\begin{array}{cc}
\inc{XN}&\inc{YN}\\
X_N&Y_N
\end{array}
\]
\caption{理论 $X_N$ 与理论 $Y_N$。\label{fig:XNYN}}
\end{figure}

为这些理论引入名称。$X_N$ 理论是对应于一个带有一个正则穿孔和一个 $N$ 阶极点穿孔的球面的超共形场论；$Y_N$ 理论是对应于一个仅带有一个 $N$ 阶极点穿孔的球面的超共形场论，见图 \ref{fig:XNYN}。从图 \ref{fig:H1}、图 \ref{fig:H2} 和图 \ref{fig:H3} 可以看出，
\begin{equation}
Y_7=AD_{N_f=1}(\SU(2)),\qquad
Y_8=AD_{N_f=2}(\SU(2))=X_5,\quad
AD_{N_f=3}(\SU(2))=X_6.\label{accidents}
\end{equation}

另外，回顾图 \ref{fig:lessG} 中给出的、含有一个味的 $\SU(2)$ 理论的构造。其中，带有一个正则穿孔和一个三阶极点穿孔的球面充当空的边界条件；带有一个正则穿孔和一个四阶极点穿孔的球面，则表现为一个处于 $\SU(2)$ 二重态中的自由超多重态。等价地说，
\begin{equation}
X_3=\text{空理论},\qquad
X_4=\text{处于 $\SU(2)$ 二重态中的自由超多重态}.
\end{equation}
图 \ref{fig:wild} 的第一排展示了它们。

\begin{figure}[h]
\[
\inc{wildtheoriesXX}
\]
\caption{通常的规范理论与野性规范理论。\label{fig:wild}}
\end{figure}
更一般地，可以考虑带有任意阶数 $N$ 和 $N'$ 极点的二穿孔球面。图 \ref{fig:wild} 第二排展示了一个 $N=6$、$N'=5$ 的例子。它可以理解为一个物质内容略显不同寻常的 $\SU(2)$ 规范理论，其物质由两个强相互作用共形场论 $X_N$ 和 $X_{N'}$ 描述。注意，二阶极点总是带有一个 $\SU(2)$ 味对称性，因此 $X_N$ 理论总是具有一个 $\SU(2)$ 味对称性。来自管的 $\SU(2)$ 规范对称性将这两个理论耦合起来。这类将 $X_N$ 作为部分物质内容的规范理论，通常称为野性规范理论。

这一理论的耦合跑动很容易求得。与往常一样，假设 $a$ 很大。$\lambda^2=\phi(z)$ 的分支点位于
\begin{equation}
\frac{\Lambda^2}{z^N}dz^2\sim\frac{udz^2}{z^2}\quad\text{或}\quad
\Lambda^2z^{N'}dz^2\sim\frac{udz^2}{z^2}.
\end{equation}
成立的位置附近。因此，它们位于
\begin{equation}
z_-\sim\left(\frac\Lambda a\right)^{2/(N-2)},\qquad
z_+\sim\left(\frac a\Lambda\right)^{2/(N'-2)}.
\end{equation}
附近。我们得到
\begin{equation}
a_D\sim\frac2{2\pi i}\int_{z_+}^{z_-}x\frac{dz}z
\sim-\frac2{2\pi i}(\frac2{N-2}+\frac2{N'-2})a\log\frac a\Lambda.
\end{equation}
这意味着，单圈跑动为
\begin{equation}
\Lambda\frac d{d\Lambda}\tau=\frac2{2\pi i}(b_N+b_{N'}-4)
\end{equation}
其中
\begin{equation}
b_N=2-\frac2{N-2}.
\end{equation}

\begin{figure}[h]
\[
\inc{cc}
\]
\caption{耦合的跑动衡量流的两点关联函数。\label{fig:cc}}
\end{figure}

一个二重态超多重态对单圈跑动的贡献为 $b=1$。因此，可以粗略地将 $b_N$ 理解为理论 $X_N$ 所具有的二重态超多重态的有效数目。更精确地说，$b_N$ 衡量的是 $\SU(2)$ 味对称性之对称流 $j_\mu$ 的关联函数的系数，见图 \ref{fig:cc}。如图所示，对于含有味的 $\SU(2)$ 理论，规范耦合的跑动由规范多重态的圈（用波浪线表示）或超多重态的圈（用直线表示）引起，它们通过 $\SU(2)$ 流算符 $j_\mu$ 与规范场耦合。于是，对单圈跑动的贡献衡量 $\vev{j_\mu j_\nu}$。$X_N$ 理论的贡献是一个自由味的 $b_N$ 倍，这就意味着
\begin{equation}
\vev{j_\mu j_\nu}_{X_N}=b_N\vev{j_\mu j_\nu}_{\text{$\SU(2)$ 二重态中的自由超多重态}}.\label{GOO}
\end{equation}
% [Tutor Note: p.122 原文 "the $X_N$ theory contributes $b_N$ times a free flavor does means that" 中 "does means" 应为 "does mean"（或直接作 "means"），属原文语法笔误；中文语义不受影响，按 "这就意味着" 译出。]
回顾 $X_3$ 只是空理论，而 $X_4$ 是一个处于味 $\SU(2)$ 二重态中的自由超多重态。我们的一般公式正确地重现了 $b_3=0$ 和 $b_4=1$。

下一节将看到，纯 $\SU(N)$ 规范理论的一个奇异极限变成理论 $Y_{N+4}$，而纯 $\SO(2N)$ 规范理论的一个奇异极限变成理论 $X_{N+2}$。我们还将看到，含有两个味的 $\SU(N)$ 规范理论具有一个由 $X_{N+3}$ 给出的奇异极限。

将由纯 $G$ 规范理论得到的 Argyres-Douglas 共形场论记为 $AD_{N_f=0}(G)$，将由含有两个味的 $\SU(N)$ 理论得到的 Argyres-Douglas 共形场论记为 $AD_{N_f=2}(\SU(N))$。于是，可以将这些等价关系简洁地写成
\begin{equation}
\begin{array}{r@{}ll}
AD_{N_f=0}&(\SU(N))&=Y_{N+4},\\
AD_{N_f=0}&(\SO(2N))&=X_{N+2},\\
AD_{N_f=2}&(\SU(N))&=X_{N+3}.
\end{array}\label{ppp}
\end{equation}

在式 \eqref{accidents} 中已经看到，来自含有 $N_f=2$ 个味的 $\SU(2)$ 理论的 Argyres-Douglas 理论 $AD_{N_f=2}(\SU(2))$，既等价于 $X_5$，也等价于 $Y_8$。从式 \eqref{ppp} 的观点来看，这一巧合是等价关系 $\SU(4)\simeq\SO(6)$ 的体现。另外，考虑纯 $\SO(4)$ 理论，它是纯 $\SU(2)$ 理论的两个副本。它最奇异的点，就是两个副本都位于单极子点的地方，从而实现两个自由超多重态。确实，这一系统具有 $\SU(2)$ 味对称性，并在其下构成二重态，实现了如下事实：
\begin{equation}
AD_{N_f=0}(\SO(4))=X_4=\text{处于 $\SU(2)$ 二重态中的一个自由超多重态}.
\end{equation}
作者觉得这些等式相当有趣，因为它们是量子场论之间的等式，而不是单个量子场论中物理量之间的等式。
